% Options for packages loaded elsewhere
\PassOptionsToPackage{unicode}{hyperref}
\PassOptionsToPackage{hyphens}{url}
\PassOptionsToPackage{dvipsnames,svgnames,x11names}{xcolor}
\documentclass[
  10pt,
  fontset=fandol]{ctexart}
\usepackage{xcolor}
\usepackage[margin=16mm]{geometry}
\usepackage{amsmath,amssymb}
\setcounter{secnumdepth}{-\maxdimen} % remove section numbering
\usepackage{iftex}
\ifPDFTeX
  \usepackage[T1]{fontenc}
  \usepackage[utf8]{inputenc}
  \usepackage{textcomp} % provide euro and other symbols
\else % if luatex or xetex
  \usepackage{unicode-math} % this also loads fontspec
  \defaultfontfeatures{Scale=MatchLowercase}
  \defaultfontfeatures[\rmfamily]{Ligatures=TeX,Scale=1}
\fi
\usepackage{lmodern}
\ifPDFTeX\else
  % xetex/luatex font selection
\fi
% Use upquote if available, for straight quotes in verbatim environments
\IfFileExists{upquote.sty}{\usepackage{upquote}}{}
\IfFileExists{microtype.sty}{% use microtype if available
  \usepackage[]{microtype}
  \UseMicrotypeSet[protrusion]{basicmath} % disable protrusion for tt fonts
}{}
\makeatletter
\@ifundefined{KOMAClassName}{% if non-KOMA class
  \IfFileExists{parskip.sty}{%
    \usepackage{parskip}
  }{% else
    \setlength{\parindent}{0pt}
    \setlength{\parskip}{6pt plus 2pt minus 1pt}}
}{% if KOMA class
  \KOMAoptions{parskip=half}}
\makeatother
\setlength{\emergencystretch}{3em} % prevent overfull lines
\providecommand{\tightlist}{%
  \setlength{\itemsep}{0pt}\setlength{\parskip}{0pt}}
\usepackage{amsmath,amssymb,mathtools,needspace}
\xeCJKDeclareCharClass{CJK}{"2032,"0400->"04FF,"2160->"216B,"2460->"2473,"25A0,"25C7,"25CB,"25CF}
\IfFontExistsTF{Arial Unicode MS}{\xeCJKsetup{AutoFallBack=true}\setCJKfallbackfamilyfont{\CJKrmdefault}{Arial Unicode MS}}{}
\setcounter{secnumdepth}{0}
\setlength{\emergencystretch}{2em}
\allowdisplaybreaks
\usepackage{bookmark}
\IfFileExists{xurl.sty}{\usepackage{xurl}}{} % add URL line breaks if available
\urlstyle{same}
\hypersetup{
  pdftitle={要避免两相近数相减},
  colorlinks=true,
  linkcolor={Maroon},
  filecolor={Maroon},
  citecolor={Blue},
  urlcolor={Blue},
  pdfcreator={LaTeX via pandoc}}

\title{要避免两相近数相减}
\author{}
\date{}

\begin{document}
\maketitle

\paragraph{1.4.2 要避免两相近数相减}\label{ux8981ux907fux514dux4e24ux76f8ux8fd1ux6570ux76f8ux51cf}

在数值计算中，两相近数相减会导致有效数字严重损失.例如，\(x=532.65, y=532.52\)都有五位有效数字，但 \(x-y=0.13\)只有两位有效数字.必须尽量避免出现这类运算，最好是改变计算方法，防止这种现象产生.现举例说明.

\textbf{例1.7} 计算 \(A=10^7(1-\cos 2^\circ)\).

\textbf{解} 由于 \(\cos 2^\circ=0.9994\)，直接计算，得

\[
A=10^7(1-\cos 2^\circ)=10^7(1-0.9994)=6 \times 10^3,
\]

只有一位有效数字；若利用 \(1-\cos x=2\sin^2\frac{x}{2}\) 计算，则有

\[
A=10^7(1-\cos 2^\circ)=2 \times (\sin 1^\circ)^2 \times 10^7=6.13 \times 10^3,
\]

具有三位有效数字（这里 \(\sin 1^\circ=0.0175\)）.

例1.7 说明，可通过改变计算公式避免或减少有效数字的损失.类似地，如果 \(x_1\) 和 \(x_2\) 很接近，则

\[
\lg x_1-\lg x_2=\lg \frac{x_1}{x_2}.
\]

用右端算式，有效数字就不会损失.当 \(x\) 很大时，有

\[
\sqrt{x+1}-\sqrt{x}=\frac{1}{\sqrt{x+1}+\sqrt{x}},
\]

都用右端算式代替左端.一般情况下，当 \(f(x)\approx f(x^*)\) 时，可用 Taylor 展开

\[
f(x)-f(x^*)=f'(x^*) (x-x^*)+\frac{f''(x^*)}{2} (x-x^*)^2+\cdots,
\]

\begin{quote}
校注（agent 补充）：本页末尾 Taylor 展开后的讨论续下页。
\end{quote}

\begin{quote}
校注（agent 补充，CH01-E005）：例1.6续文原书写``用小数除以大数''，并在第二种消去结果的首行两侧印 \(10^6\)，此处均照录。按上下文，第一种操作是将第一方程除以小量 \((10^{-4}\times0.1000)/2\)；第二种操作直接以 \(R_1-0.000005R_2\) 消去，四位浮点结果应得到 \(10^1\times0.1000x_2=10^1\times0.1000\)。原印 \(10^6\) 的等式虽可通过两侧再同乘 \(10^5\) 得到，但原文没有该缩放步骤；上述措辞和指数均列为疑误，未静默修改。
\end{quote}

\begin{quote}
校注（agent 补充，CH01-E006）：例1.7原文``\(6.13\times10^3\)，具有三位有效数字''照录。独立数值复算 \(10^7(1-\cos2^\circ)\approx6091.729809\)，原列 \(6130\) 的误差约 \(38.270191\)，超过三位有效数字所对应的 \(5\)。以本章有效数字定义衡量，该``三位''判断存在问题；这不影响使用半角公式减轻相消的思路。
\end{quote}

\href{../../source-images/pdf-024.jpeg}{查看原页}

取右端的有限项近似左端. 如果无法改变算式, 则采用增加有效位数进行运算; 在计算机上则采用双倍字长运算, 但这要增加机器计算时间和多占内存单元.

\begin{center}\rule{0.5\linewidth}{0.5pt}\end{center}

\subsection{相关原页校注（agent 补充）}\label{ux76f8ux5173ux539fux9875ux6821ux6ce8agent-ux8865ux5145}

以下校注来自本节覆盖的原页，可能也涉及同页相邻小节；原文照录与校正说明分开保留。

\subsubsection{原书 PDF 页 24 的校注}\label{ux539fux4e66-pdf-ux9875-24-ux7684ux6821ux6ce8}

\href{../pages/pdf-024.md}{完整原页转录} · \href{../../source-images/pdf-024.jpeg}{原页图像}

\begin{quote}
校注（agent 补充）：本页末尾乘法次数公式的说明续下页``次乘法和 \(n\) 次加法''。
\end{quote}

\subsection{本节覆盖原页的校注记录}\label{ux672cux8282ux8986ux76d6ux539fux9875ux7684ux6821ux6ce8ux8bb0ux5f55}

下列记录按原页关联，含同页相邻小节的校注；计算时同时读取上文校注和完整原页。

\begin{itemize}
\tightlist
\item
  \texttt{CH01-E005}：原书 PDF 页 23
\item
  \texttt{CH01-E006}：原书 PDF 页 23
\end{itemize}

\end{document}
